\kefalaio{Φυσικοί αριθμοί: η σειρά που δεν τελειώνει}{Κεφάλαιο Α.1.1, «Φυσικοί αριθμοί – Διάταξη φυσικών – Στρογγυλοποίηση»}

\begin{anagke}
Δεύτερη μέρα στο γυμνάσιο. Στην είσοδο μιας αίθουσας ένα χαρτί λέει: «Το σχολείο μας έχει 387 μαθητές». Ο Άρης το φωτογραφίζει και στέλνει στον αδελφό του τον Θοδωρή: «Το σχολείο έχει 400 παιδιά!». Ο Θοδωρής απαντάει με ένα χαμογελαστό «στρογγύλεψες».

Στο μάθημα η κυρία Αλεξάνδρα δεν ανοίγει βιβλίο. Γράφει στον πίνακα τρεις αριθμούς: τους μαθητές του σχολείου, 387, το υψόμετρο του Ολύμπου, 2.917 μέτρα, και τον αριθμό του τηλεφώνου της γραμματείας. «Σας ζητάω να τους πείτε σε ένα φίλο σας. Ποιον θα πείτε ακριβώς και ποιον περίπου;»

«Όλους ακριβώς», λέει ο Στέλιος. «Τι νόημα έχει το “περίπου”;». Η Ζωή σηκώνει το χέρι. «Αν στον φίλο σου πεις ότι περπάτησες 8.412 βήματα, θα θυμάται μόνο το 8 και θα μπερδεύεται με τα υπόλοιπα. Το “περίπου 8.000” το κρατάς». «Και πώς ξέρω ότι όλοι θα βγάλουμε το ίδιο “περίπου”;» ρωτάει η Λίνα. «Αν ο Άρης κάνει το 387 τετρακόσια και εγώ τριακόσια ογδόντα;»

Η κυρία Αλεξάνδρα χαμογελάει. Η ερώτηση της Λίνας είναι ακριβώς αυτή που ήθελε να ακούσει. Πότε επιτρέπεται το «περίπου», και πώς γίνεται με κανόνα, ώστε όλοι να κάνουμε το ίδιο;
\end{anagke}

\begin{lexi}{φυσικός, στρογγυλοποίηση}{φυσικός, στρογγυλοποίηση}
Οι \lexiword{φυσικοί} αριθμοί λέγονται έτσι επειδή είναι οι πρώτοι που βρίσκει κανείς στη φύση, όταν μετράει πράγματα: μήλα, παιδιά, βήματα. Κανείς δεν μέτρησε ποτέ «ενάμισι παιδί». Το μηδέν μπήκε στην παρέα αργότερα. Στο σχολείο θα το βρεις μέσα στους φυσικούς.

Η \lexiword{στρογγυλοποίηση} είναι το να κάνεις έναν αριθμό «στρογγυλό», δηλαδή να τον αλλάξεις με έναν κοντινό του που τελειώνει σε μηδενικά: 387 γίνεται 400. Στρογγυλός σημαίνει τακτοποιημένος, χωρίς αιχμές. Δεν είναι μαγικό: είναι \emph{συμφωνία}. Όλοι κάνουμε τα ίδια βήματα και βγάζουμε το ίδιο αποτέλεσμα.

Δύο ακόμα λέξεις που θα ξανασυναντήσεις: \lexiword{άρτιος} είναι ο αριθμός που «ταιριάζει» τέλεια σε ζευγάρια, και \lexiword{περιττός} αυτός που «περισσεύει» ένας. Τους λέμε και ζυγούς και μονούς.
\end{lexi}

\begin{metafrasi}
Ο αριθμός που ακολουθεί το 9. & $9+1=10$\\
Το πέντε είναι μεγαλύτερο από το τρία. & $5>3$\\
Ο αριθμός $x$ βρίσκεται ανάμεσα στο 3 και στο 8. & $3<x<8$\\
Τρία εκατομμύρια, σαράντα χιλιάδες και επτά. & 3.040.007\\
Το 9.573.842 στρογγυλοποιημένο στις εκατοντάδες. & $9.573.842\approx 9.573.800$\\
Ο αριθμός είναι άρτιος. & Διαιρείται με το 2\\
\end{metafrasi}

\tmima{Η θεωρία}

\begin{orismos}{φυσικοί αριθμοί}
Οι αριθμοί $0,\,1,\,2,\,3,\,4,\,\dots,\,99,\,100,\,\dots$ λέγονται \textbf{φυσικοί αριθμοί}. Κάθε φυσικός έχει έναν επόμενο και, εκτός από το 0, έναν προηγούμενο.
\end{orismos}

Οι φυσικοί χωρίζονται σε δύο κατηγορίες: \textbf{άρτιους} (ή ζυγούς), που διαιρούνται με το 2, και \textbf{περιττούς} (ή μονούς), που δεν διαιρούνται. Ο άρτιος τελειώνει πάντα σε 0, 2, 4, 6 ή 8.

\tmimasmall{Δέκα ψηφία, άπειροι αριθμοί}

Με τα δέκα ψηφία $0,\,1,2,\dots,9$ γράφουμε όσους αριθμούς θέλουμε. Αυτό γίνεται επειδή η αξία ενός ψηφίου εξαρτάται και από τη \textbf{θέση} του, δηλαδή από τη \textbf{δεκαδική τάξη}: μονάδες, δεκάδες, εκατοντάδες, χιλιάδες, δεκάδες χιλιάδων, εκατοντάδες χιλιάδων, εκατομμύρια. Το 0 κρατάει μια θέση άδεια: στο 5.780.901 τα δύο μηδενικά λένε ότι δεν υπάρχουν ούτε χιλιάδες ούτε δεκάδες.

\tmimasmall{Διάταξη: ποιος είναι μεγαλύτερος}

Γράφουμε $=$ για «ίσος με», $<$ για «μικρότερος από» και $>$ για «μεγαλύτερος από». Δύο φυσικούς μπορούμε πάντα να τους συγκρίνουμε.

\begin{kanonas}{πώς συγκρίνω δύο φυσικούς}
\begin{enumerate}[leftmargin=1.6em,itemsep=0.2ex,topsep=2pt]
\item Μετράω τα ψηφία. Ο αριθμός με περισσότερα ψηφία είναι μεγαλύτερος (αν δεν ξεκινάει με 0).
\item Αν έχουν ίδιο πλήθος ψηφίων, τα κοιτάζω από αριστερά προς τα δεξιά. Στο πρώτο ψηφίο που διαφέρουν, μεγαλύτερος είναι αυτός με το μεγαλύτερο ψηφίο.
\end{enumerate}
\end{kanonas}

Έτσι $90.876>86.945$, γιατί στο πρώτο ψηφίο το 9 είναι μεγαλύτερο από το 8, και $345<5.690$, γιατί το 345 έχει λιγότερα ψηφία. Τους φυσικούς μπορούμε να τους βάλουμε σε σειρά: \textbf{αύξουσα}, από τον μικρότερο στον μεγαλύτερο, ή \textbf{φθίνουσα}, αντίστροφα. Δεν υπάρχει μεγαλύτερος φυσικός αριθμός, γιατί κάθε φυσικός έχει επόμενο. Ο μικρότερος είναι το 0.

\tmimasmall{Η αριθμητική ευθεία}

Η διάταξη μας επιτρέπει να βάλουμε τους φυσικούς πάνω σε μια ευθεία. Διαλέγουμε ένα σημείο Ο για το 0 (την αρχή) και ένα σημείο Α δεξιά του για το 1. Το τμήμα ΟΑ είναι η \textbf{μονάδα μέτρησης}. Με το ίδιο μήκος βρίσκουμε τα σημεία των 2, 3, 4, \dots

\begin{center}
\begin{tikzpicture}[x=1.15cm,y=1cm]
\draw[-{Stealth}] (-0.3,0)--(9.6,0);
\foreach \i in {0,...,8}{\draw (\i,-0.12)--(\i,0.12); \node[below,font=\small] at (\i,-0.12) {\i};}
\node[above,font=\small\bfseries,gblue] at (0,0.15) {Ο};
\node[above,font=\small\bfseries,gblue] at (1,0.15) {Α};
\draw[<->,gblue] (0,0.75)--(1,0.75) node[midway,above,font=\scriptsize] {μονάδα};
\end{tikzpicture}
\end{center}

\tmimasmall{Στρογγυλοποίηση}

\begin{kanonas}{στρογγυλοποίηση φυσικού}
\begin{enumerate}[leftmargin=1.6em,itemsep=0.2ex,topsep=2pt]
\item Βρίσκω την τάξη στην οποία θα στρογγυλοποιήσω (δεκάδες, εκατοντάδες, χιλιάδες, \dots).
\item Κοιτάζω το ψηφίο της \textbf{αμέσως μικρότερης} τάξης.
\item Αν είναι $0,\,1,2,\,3$ ή $4$, το ψηφίο της τάξης μου μένει ίδιο. Αν είναι $5,\,6,7,\,8$ ή $9$, αυξάνεται κατά 1.
\item Όλα τα ψηφία δεξιά από την τάξη μου γίνονται 0.
\end{enumerate}
\end{kanonas}

Το αποτέλεσμα λέγεται \textbf{προσέγγιση} του αριθμού και γράφεται με το σύμβολο $\approx$ (διαβάζεται «περίπου ίσο»). Για παράδειγμα, $2.917\approx 3.000$ στις χιλιάδες και $2.917\approx 2.900$ στις εκατοντάδες.

Η στρογγυλοποίηση ταιριάζει εκεί που μας ενδιαφέρει η τάξη μεγέθους, όπως το πλήθος των ανθρώπων σε μια πόλη. Δεν ταιριάζει εκεί που χρειάζεται ακρίβεια: ο αριθμός του τηλεφώνου, το ΑΦΜ, ο κωδικός ενός λουκέτου.

\begin{prosoxi}
\begin{itemize}[leftmargin=1.2em,itemsep=0.3ex,topsep=2pt]
\item Στρογγυλοποιούμε \textbf{μία φορά}, από τον αρχικό αριθμό. Το 1.449 στις εκατοντάδες είναι 1.400, γιατί το ψηφίο των δεκάδων είναι 4. Αν το κάνεις σε δύο βήματα, πρώτα στις δεκάδες (1.450) και μετά στις εκατοντάδες, βγάζεις 1.500, που είναι λάθος.
\item Όταν το ψηφίο που αυξάνεται είναι 9, η αύξηση «κουβαλάει»: $9.995\approx 10.000$ στις δεκάδες.
\item Όταν μετράς \textbf{από ως}, και τα δύο άκρα μετράνε. Από τη σελίδα 12 ως τη σελίδα 35 είναι $35-12+1=24$ σελίδες, όχι 23.
\end{itemize}
\end{prosoxi}

\tmima{Λυμένα παραδείγματα}

\paradeigma{Οι κωδικοί του ντουλαπιού}
\tixero Το ντουλάπι του Ηλία ανοίγει με έναν τριψήφιο κωδικό. Θυμάται ότι έχει τα ψηφία 4, 0 και 7, καθένα μία φορά, αλλά όχι με ποια σειρά.\par
\tithelo Όλους τους πιθανούς κωδικούς, σε αύξουσα σειρά.\par
\possindeo Δοκιμάζω όλες τις σειρές των τριών ψηφίων: 407, 470, 047, 074, 704, 740. Ο 047 και ο 074 δεν είναι τριψήφιοι αριθμοί, γιατί το 0 μπροστά δεν μετράει. Τους σβήνω.\par
\apantisi 407, 470, 704, 740.\par
\elegxos Είναι τέσσερις, και κανένας δεν ξεκινάει από 0. Ο μικρότερος είναι ο 407, γιατί ξεκινάει από το μικρότερο μη μηδενικό ψηφίο, το 4, και μετά βάζει το μικρότερο που μένει, το 0.

\paradeigma{Στρογγυλοποιώ το ίδιο νούμερο τρεις φορές}
\tixero Ο αριθμός 9.573.842.\par
\tithelo Να τον στρογγυλοποιήσω (α) στις εκατοντάδες, (β) στις χιλιάδες, (γ) στα εκατομμύρια.\par
\possindeo
(α) Τάξη: εκατοντάδες. Η αμέσως μικρότερη τάξη είναι οι δεκάδες, με ψηφίο 4. Επειδή $4<5$, το 8 μένει ίδιο και τα δεξιά γίνονται 0: $9.573.842\approx 9.573.800$.\par
(β) Τάξη: χιλιάδες, ψηφίο 3. Η αμέσως μικρότερη τάξη, οι εκατοντάδες, έχει ψηφίο 8. Επειδή $8\ge 5$, το 3 γίνεται 4: $9.573.842\approx 9.574.000$.\par
(γ) Τάξη: εκατομμύρια, ψηφίο 9. Η εκατοντάδα χιλιάδων έχει ψηφίο 5. Επειδή $5\ge 5$, το 9 αυξάνεται και γίνεται 10: $9.573.842\approx 10.000.000$.\par
\apantisi 9.573.800, 9.574.000 και 10.000.000.\par
\elegxos Κάθε προσέγγιση απέχει λιγότερο από μισή τάξη από τον αρχικό. Για παράδειγμα, στο (β), το 9.573.842 απέχει 158 από το 9.574.000 και 842 από το 9.573.000. Το πρώτο είναι πιο κοντά.

\paradeigma{Πόσες σελίδες;}
\tixero Η Φανή διαβάζει από τη σελίδα 47 ως τη σελίδα 112 του βιβλίου της, και τις δύο μαζί.\par
\tithelo Πόσες σελίδες διαβάζει.\par
\possindeo Αν έβγαζα $112-47=65$, θα είχα αφήσει έξω την πρώτη σελίδα. Για να μετρήσω και τα δύο άκρα, προσθέτω 1: $112-47+1=66$.\par
\apantisi 66 σελίδες.\par
\elegxos Δοκιμή με μικρούς αριθμούς: από τη σελίδα 3 ως τη 5 είναι τρεις σελίδες (3, 4, 5). Ο τύπος δίνει $5-3+1=3$. Σωστά.

\tmima{Ασκήσεις}

\zoni{Το κατάλαβα}{εφαρμογή του κανόνα}

\ex Γράψε με ψηφία: α) διακόσια πέντε \quad β) επτακόσια τριάντα δύο \quad γ) είκοσι χιλιάδες οκτακόσια δεκατρία \quad δ) τέσσερα εκατομμύρια πεντακόσιες χιλιάδες δεκαπέντε.

\ex Γράψε με λέξεις: α) 38.951 \quad β) 5.000.812 \quad γ) 120.003.

\ex Ποιοι είναι οι τρεις προηγούμενοι αριθμοί του 289 και ποιοι οι δύο επόμενοι; Ποιος είναι ο προηγούμενος του 1.000.000 και ποιος ο προηγούμενος του 0;

\ex Γράψε σε αύξουσα σειρά τους αριθμούς 3.515, 4.800, 3.620, 3.508 και 4.801. Μετά γράψε τους ίδιους σε φθίνουσα.

\ex Βάλε το σωστό σύμβολο ($<$, $=$ ή $>$): α) $45\ \square\ 45$ \quad β) $38\ \square\ 36$ \quad γ) $456\ \square\ 465$ \quad δ) $8.765\ \square\ 8.970$ \quad ε) $90.876\ \square\ 86.945$ \quad στ) $345\ \square\ 5.690$.

\ex Σχεδίασε μια αριθμητική ευθεία με αρχή το Ο και μονάδα $ΟΑ=2$ εκ. Τοποθέτησε τα σημεία Β, Γ, Δ, Ε σε αποστάσεις 6, 10, 12 και 14 εκ. από το Ο. Ποιοι αριθμοί αντιστοιχούν στα σημεία αυτά;

\ex Στρογγυλοποίησε στην πλησιέστερη εκατοντάδα: 345, 761, 659, 2.567, 9.532, 31.549 και 8.950.

\ex Στρογγυλοποίησε τον αριθμό 7.568.349 (α) στις δεκάδες, (β) στις εκατοντάδες, (γ) στις χιλιάδες, (δ) στις δεκάδες χιλιάδων, (ε) στα εκατομμύρια.

\ex Σωστό ή λάθος;\par
α) Στον αριθμό 5.780.901 τα δύο μηδενικά δείχνουν ότι δεν υπάρχουν χιλιάδες και δεκάδες.\quad β) Δέκα χιλιάδες είναι μία δεκάδα χιλιάδων.\quad γ) Από τον 32 ως τον 122 υπάρχουν 91 αριθμοί.\quad δ) Από τη σελίδα 12 ως τη σελίδα 35 είναι 24 σελίδες.\quad ε) Σε οκτώ ημέρες από σήμερα, που είναι Πέμπτη, θα είναι Παρασκευή.\quad στ) Ανάμεσα στο 2 και στο 3 δεν υπάρχει άλλος φυσικός αριθμός.\quad ζ) Το 1.449 στρογγυλοποιημένο στις εκατοντάδες είναι 1.500.

\zoni{Το χρησιμοποιώ}{προβλήματα από τη ζωή}

\ex \textbf{Το ντουλάπι της Ζωής.} Ο κωδικός του ντουλαπιού της Ζωής είναι τριψήφιος αριθμός με τα ψηφία 3, 0 και 6, καθένα μία φορά. Γράψε όλους τους δυνατούς κωδικούς από τον μικρότερο στον μεγαλύτερο. Γιατί δεν υπάρχει ο κωδικός 036;

\ex \textbf{Ο Όλυμπος.} Η κορυφή του Ολύμπου, ο Μύτικας, έχει υψόμετρο 2.917 μέτρα. Η Λίνα γράφει σε μια πινακίδα για το σχολικό ημερολόγιο «περίπου 3.000 μ.» και ο Ηλίας «περίπου 2.900 μ.». Ποιος στρογγυλοποίησε σε χιλιάδες και ποιος σε εκατοντάδες; Ποιος από τους δύο λέει κάτι πιο κοντά στην αλήθεια;

\ex \textbf{Τα βήματα του Άρη.} Η εφαρμογή στο κινητό του Άρη γράφει ότι σήμερα έκανε 9.987 βήματα. Ο Άρης λέει «Στρογγυλά, 10.000! Πέτυχα τον στόχο μου». Ο στόχος είναι τουλάχιστον 10.000 βήματα. Έφτασε τον στόχο; Πόσα βήματα του λείπουν; Γιατί εδώ η στρογγυλοποίηση δεν βοηθάει;

\ex \textbf{Τα πλακάκια του θείου.} Ο θείος Λευτέρης χρειάζεται 1.234 πλακάκια για τη βεράντα του. Πωλούνται σε κιβώτια των 100. Ο Στέλιος στρογγυλοποιεί στην πλησιέστερη εκατοντάδα και λέει «12 κιβώτια». Φτάνουν; Πόσα κιβώτια πρέπει να αγοράσει; Πόσα πλακάκια θα περισσέψουν;

\ex \textbf{Το βιβλίο της Μυρτώς.} Η Μυρτώ διαβάζει από τη σελίδα 47 ως τη σελίδα 112, και τις δύο μαζί. Πόσες σελίδες διαβάζει; Αν διαβάζει 6 σελίδες την ημέρα, φτάνουν 10 ημέρες;

\ex \textbf{Ο φράχτης.} Ο παππούς φράζει το χωράφι του κατά μήκος ενός δρόμου 20 μ. Βάζει έναν πάσσαλο κάθε 2 μ., και έναν στην αρχή και έναν στο τέλος. Πόσους πασσάλους θα βάλει;

\ex \textbf{Τα εισιτήρια.} Σε μια παράσταση τα εισιτήρια είναι αριθμημένα από το 120 ως το 399. Πόσα εισιτήρια υπάρχουν; Πόσα από αυτά έχουν περιττό αριθμό;

\ex \textbf{Ο άξονας του Ηλία.} Ο Ηλίας σχεδιάζει έναν άξονα που δείχνει το 0 και το 100 (δες το σχήμα). Ποιοι αριθμοί αντιστοιχούν στα σημεία Α, Β, Γ και Δ;

\begin{center}
\begin{tikzpicture}[x=1cm]
\draw[-{Stealth}] (-0.3,0)--(11.2,0);
\draw (0,-0.12)--(0,0.12) node[above,font=\small] {0};
\draw (2,-0.12)--(2,0.12) node[above,font=\small] {100};
\foreach \x/\n in {3/Α,5/Β,7.5/Γ,10/Δ}{\draw[gblue,thick] (\x,-0.15)--(\x,0.15); \node[below,font=\small\bfseries,gblue] at (\x,-0.15) {\n};}
\end{tikzpicture}
\end{center}

\zoni{Το σκέφτομαι}{μια ερώτηση που πάει πιο μακριά}

\ex Πόσες φορές γράφεται το ψηφίο 9, αν γράψεις όλους τους αριθμούς από το 1 ως το 100;

\ex Ποιος είναι ο μικρότερος και ποιος ο μεγαλύτερος φυσικός αριθμός που, στρογγυλοποιημένος στις χιλιάδες, δίνει 3.000;

\ex Ποιος είναι ο μεγαλύτερος και ποιος ο μικρότερος τετραψήφιος αριθμός που έχει όλα τα ψηφία του διαφορετικά;

\ex Ο Άρης λέει: «Θα βρω τον μεγαλύτερο φυσικό αριθμό και θα τον γράψω στον πίνακα». Η Ζωή λέει: «Δεν υπάρχει». Εξήγησε με λόγια γιατί έχει δίκιο η Ζωή.

\ex Ο Ηλίας δοκιμάζει στο κινητό: $7+5=12$, $3+9=12$. Ρωτάει τη Λίνα: «Αν προσθέσω δύο φυσικούς, βγαίνει πάντα φυσικός. Αν αφαιρέσω;» Η Λίνα γράφει $3-5$. Μπορεί να υπάρχει αυτό το αποτέλεσμα στους φυσικούς; Τι νομίζεις ότι θα χρειαστούμε;
